\documentclass[11pt]{article}

\usepackage[T1]{fontenc}
\usepackage[utf8]{inputenc}
\usepackage{lmodern}
\usepackage{amsmath,amssymb,amsthm,mathtools}
\usepackage{microtype}
\usepackage{booktabs}
\usepackage{enumitem}
\usepackage{tabularx}
\usepackage{needspace}
\usepackage{array}
\usepackage[margin=1in,top=0.9in,bottom=1in]{geometry}
\usepackage{fancyhdr}
\usepackage{lastpage}
\usepackage{tikz}
\usepackage[hidelinks]{hyperref}

\pagestyle{fancy}
\fancyhf{}
\lhead{\small BIO 2970 --- Selection Practice Exam S3 Key}
\rhead{\small Answer Key}
\cfoot{\small\thepage\ of \pageref*{LastPage}}
\renewcommand{\headrulewidth}{0.4pt}
\setlength{\headheight}{14pt}

% answer text, for KEY documents only
\newcommand{\keyans}[1]{\par\medskip\noindent\textit{Answer.}\ #1\par\medskip}
\newcommand{\work}[1]{\par\vspace{#1}}
\newcommand{\opt}[1]{\makebox[1.9em][l]{(#1)}}
\newcolumntype{Y}{>{\raggedright\arraybackslash}X}
\setlist[enumerate,1]{leftmargin=2em,labelsep=0.6em}
\hyphenpenalty=1000
\relpenalty=10000
\binoppenalty=10000
\begin{document}

\begin{center}
  {\LARGE\bfseries Natural Selection --- Practice Exam S3 \textit{Answer Key}}\\[5pt]
  {\large Level 3 \quad\textbullet\quad Hardest}\\[6pt]
  BIO 2970 \quad\textbullet\quad Fitness and Selection, Lectures 23 and 24 \quad\textbullet\quad Fall 2026\\[3pt]
  \rule{\textwidth}{0.8pt}\\[3pt]
  {\small Every answer is explained, including why the tempting wrong choices fail. Grade the answer sheet with the grid below, then read the explanations for anything you missed.}
\end{center}
\vspace{0.4em}

\begin{center}\small\setlength{\tabcolsep}{6.2pt}\begin{tabular}{lcccccccccccc}
\toprule
Question & 1 & 2 & 3 & 4 & 5 & 6 & 7 & 8 & 9 & 10 & 11 & 12 \\
\midrule
Answer & \textbf{D} & \textbf{C} & \textbf{D} & \textbf{A} & \textbf{B} & \textbf{B} & \textbf{C} & \textbf{B} & \textbf{C} & \textbf{B} & \textbf{D} & \textbf{D} \\
\bottomrule\end{tabular}\end{center}
\begin{center}\small\setlength{\tabcolsep}{6.2pt}\begin{tabular}{lcccccccccccc}
\toprule
Question & 13 & 14 & 15 & 16 & 17 & 18 & 19 & 20 & 21 & 22 & 23 & 24 \\
\midrule
Answer & \textbf{A} & \textbf{D} & \textbf{A} & \textbf{B} & \textbf{C} & \textbf{C} & \textbf{A} & \textbf{A} & \textbf{D} & \textbf{B} & \textbf{A} & \textbf{C} \\
\bottomrule\end{tabular}\end{center}
\begin{center}\small\begin{tabular}{lll}\toprule Topic & Questions & Your misses\\\midrule
General equation and mean fitness & 4, 6, 10, 21 & \\
Counts or Hardy-Weinberg & 9, 18, 19 & \\
Average excess & 7, 13, 22 & \\
Types of selection & 17, 23 & \\
Special models, \(p\) and \(q\) & 2, 3, 5, 8, 16 & \\
Heterozygote advantage and \(\hat{q}\) & 1, 12, 14, 20 & \\
Heterozygote disadvantage and stability & 11 & \\
Human cases and Darwin & 15, 24 & \\
\bottomrule\end{tabular}\end{center}
\vspace{0.6em}
\begin{enumerate}[label=\textbf{\arabic*.}, leftmargin=2em, itemsep=1.4em]
\needspace{14\baselineskip}
\item Under heterozygote advantage, \(\hat{q}\) = 0.25 and the \(q\) homozygote has \(s_2\) = 0.9. What is \(s_1\)?
\nopagebreak\par\noindent{\footnotesize\itshape Topic: Heterozygote advantage and \(\hat{q}\)}
\par\smallskip\noindent\begin{tabularx}{\linewidth}{@{}YYYY@{}}
\makebox[1.9em][l]{(A)}2.70 & \makebox[1.9em][l]{(B)}0.225 & \makebox[1.9em][l]{(C)}0.25 & \makebox[1.9em][l]{\textbf{(D)}}\textbf{0.30}\ \(\checkmark\) \\
\end{tabularx}
\keyans{At equilibrium \(\hat{p} \cdot s_1\) = \(\hat{q} \cdot s_2\), so 0.75\(\cdot s_1\) = 0.25 \(\times\) 0.9 = 0.225, and \(s_1\) = 0.30. The 0.225 stops one step early, before dividing by \(\hat{p}\). The 2.70 swaps \(\hat{p}\) and \(\hat{q}\), and a selection coefficient above 1 is impossible. Check, 0.3/(0.3 + 0.9) = 0.25.}
\needspace{14\baselineskip}
\item Selection against the recessive, Hardy-Weinberg. At \(q\) = 0.4 you observe \(\Delta q\) = \(-\)0.04. What is \(s\)?
\nopagebreak\par\noindent{\footnotesize\itshape Topic: Special models, \(p\) and \(q\)}
\par\smallskip\noindent\begin{tabularx}{\linewidth}{@{}YYYY@{}}
\makebox[1.9em][l]{(A)}0.236 & \makebox[1.9em][l]{(B)}0.417 & \makebox[1.9em][l]{\textbf{(C)}}\textbf{0.391}\ \(\checkmark\) & \makebox[1.9em][l]{(D)}0.250 \\
\end{tabularx}
\keyans{Top, \(- s\)(0.6)(0.16) = \(-\)0.096s. Bottom, 1 \(-\) 0.16s. Set \(-\)0.096s/(1 \(-\) 0.16s) = \(-\)0.04 and multiply both sides by the bottom. 0.096s = 0.04 \(-\) 0.0064s, so 0.1024s = 0.04 and \(s\) = 0.391. The 0.417 forgets the bottom. The 0.250 drops \(p\), 0.04/0.16. The 0.236 uses the dominant formula.}
\needspace{14\baselineskip}
\item A lethal recessive is at \(q'\) = 0.0909 after one generation of selection. What was \(q\) before?
\nopagebreak\par\noindent{\footnotesize\itshape Topic: Special models, \(p\) and \(q\)}
\par\smallskip\noindent\begin{tabularx}{\linewidth}{@{}YYYY@{}}
\makebox[1.9em][l]{(A)}0.0909 & \makebox[1.9em][l]{(B)}0.0992 & \makebox[1.9em][l]{(C)}0.0826 & \makebox[1.9em][l]{\textbf{(D)}}\textbf{0.100}\ \(\checkmark\) \\
\end{tabularx}
\keyans{For \(s\) = 1, \(q'\) = \(q\)/(1 + \(q\)). Solve for \(q\), \(q\) = \(q'\)/(1 \(-\) \(q'\)) = 0.0909/0.9091 = 0.100. Check, 0.1/1.1 = 0.0909.}
\needspace{14\baselineskip}
\item Hardy-Weinberg, \(q\) = 0.5. AA and Aa survive fully. After selection \(q'\) = 0.4. What is the relative fitness of aa?
\nopagebreak\par\noindent{\footnotesize\itshape Topic: General equation and mean fitness}
\par\smallskip\noindent\begin{tabularx}{\linewidth}{@{}YYYY@{}}
\makebox[1.9em][l]{\textbf{(A)}}\textbf{0.333}\ \(\checkmark\) & \makebox[1.9em][l]{(B)}0.800 & \makebox[1.9em][l]{(C)}0.600 & \makebox[1.9em][l]{(D)}0.500 \\
\end{tabularx}
\keyans{\(\overline{W}\) = 0.25 + 0.5 + 0.25W. \(q'\) = (0.25 + 0.25W)/(0.75 + 0.25W) = 0.4. So 0.25 + 0.25W = 0.30 + 0.10W, 0.15W = 0.05, \(W\) = 0.333.}
\needspace{14\baselineskip}
\item Viability is 0.5, 0.5 and 0.25 and fecundity 4, 4 and 6 for AA, Aa and aa. \(q\) = 0.4, Hardy-Weinberg. What is \(\Delta q\) for a?
\nopagebreak\par\noindent{\footnotesize\itshape Topic: Special models, \(p\) and \(q\)}
\par\smallskip\noindent\begin{tabularx}{\linewidth}{@{}YYYY@{}}
\makebox[1.9em][l]{(A)}\(-\)0.043 & \makebox[1.9em][l]{\textbf{(B)}}\textbf{\(-\)0.025}\ \(\checkmark\) & \makebox[1.9em][l]{(C)}\(-\)0.052 & \makebox[1.9em][l]{(D)}\(-\)0.024 \\
\end{tabularx}
\keyans{\(w\) = 2.0, 2.0 and 1.5, so \(W\) = 1, 1 and 0.75 and \(s\) = 0.25 against the recessive. \(\Delta q\) = \(-\)0.25 \(\times\) 0.6 \(\times\) 0.16/(1 \(-\) 0.25 \(\times\) 0.16) = \(-\)0.024/0.96 = \(-\)0.025. Survival only gives \(s\) = 0.5 and \(-\)0.052. The dominant formula gives \(-\)0.043, and \(-\)0.024 forgets the bottom.}
\needspace{14\baselineskip}
\item Same population, now at \(q'\) = 0.375. What is \(\overline{W}\) in the next generation?
\nopagebreak\par\noindent{\footnotesize\itshape Topic: General equation and mean fitness}
\par\smallskip\noindent\begin{tabularx}{\linewidth}{@{}YYYY@{}}
\makebox[1.9em][l]{(A)}0.906 & \makebox[1.9em][l]{\textbf{(B)}}\textbf{0.965}\ \(\checkmark\) & \makebox[1.9em][l]{(C)}0.975 & \makebox[1.9em][l]{(D)}0.960 \\
\end{tabularx}
\keyans{Against the recessive, \(\overline{W}\) = 1 \(-\) sq\(^2\) = 1 \(-\) 0.25 \(\times\) 0.1406 = 0.965. Before selection it was 1 \(-\) 0.25 \(\times\) 0.16 = 0.960, so it rose.}
\needspace{14\baselineskip}
\item Sickle cell in Gabon, \(q\) = 0.05 for HbS, \(W\) = 0.879, 1 and 0 for HbA/HbA, HbA/HbS and HbS/HbS. What is the average excess of HbS?
\nopagebreak\par\noindent{\footnotesize\itshape Topic: Average excess}
\par\smallskip\noindent\begin{tabularx}{\linewidth}{@{}YYYY@{}}
\makebox[1.9em][l]{(A)}+0.950 & \makebox[1.9em][l]{(B)}\(-\)0.062 & \makebox[1.9em][l]{\textbf{(C)}}\textbf{+0.062}\ \(\checkmark\) & \makebox[1.9em][l]{(D)}\(-\)0.888 \\
\end{tabularx}
\keyans{\(\overline{W}\) = 0.9025(0.879) + 0.095(1) + 0 = 0.888. An HbS copy experiences 0.05(0) + 0.95(1) = 0.95. \(\alpha\) = 0.95 \(-\) 0.888 = +0.062. Positive, so rare HbS rises toward \(\hat{q}\) = 0.108.}
\needspace{14\baselineskip}
\item Same \(s\) = 0.1 and \(q\) = 0.5. Compare \(\Delta q\) for selection against the recessive and against the dominant.
\nopagebreak\par\noindent{\footnotesize\itshape Topic: Special models, \(p\) and \(q\)}
\par\smallskip
\begin{enumerate}[label={}, leftmargin=2.4em, labelwidth=0pt, itemsep=2pt, topsep=0pt]
\item \hspace*{-2.4em}\makebox[2.4em][l]{(A)}The dominant is about 50 times faster
\item \hspace*{-2.4em}\makebox[2.4em][l]{\textbf{(B)}}\textbf{About equal, \(-\)0.0128 and \(-\)0.0135}\ \(\checkmark\)
\item \hspace*{-2.4em}\makebox[2.4em][l]{(C)}The recessive is about 10 times faster
\item \hspace*{-2.4em}\makebox[2.4em][l]{(D)}They have opposite signs
\end{enumerate}
\keyans{Recessive, \(-\)0.1 \(\times\) 0.5 \(\times\) 0.25/0.975 = \(-\)0.0128. Dominant, \(-\)0.1 \(\times\) 0.25 \(\times\) 0.5/0.925 = \(-\)0.0135. When the allele is common, half its copies are already in homozygotes, so hiding barely matters. The big gap only appears when \(q\) is small.}
\needspace{14\baselineskip}
\item Researchers count 50 AA, 20 Aa and 30 aa adults, then record how many offspring each adult has. Which frequencies go into the selection calculation?
\nopagebreak\par\noindent{\footnotesize\itshape Topic: Counts or Hardy-Weinberg}
\par\smallskip
\begin{enumerate}[label={}, leftmargin=2.4em, labelwidth=0pt, itemsep=2pt, topsep=0pt]
\item \hspace*{-2.4em}\makebox[2.4em][l]{(A)}The square root of 0.3 for \(q\)
\item \hspace*{-2.4em}\makebox[2.4em][l]{(B)}Either, they give the same answer
\item \hspace*{-2.4em}\makebox[2.4em][l]{\textbf{(C)}}\textbf{The counted 0.5, 0.2 and 0.3, because selection acts on these very adults}\ \(\checkmark\)
\item \hspace*{-2.4em}\makebox[2.4em][l]{(D)}\(p^2\), \(2pq\) and \(q^2\) from \(q\) = 0.4
\end{enumerate}
\keyans{Selection here acts on the individuals counted, so use their frequencies in general equation \#1. These counts are far from Hardy-Weinberg (\(2pq\) would be 0.48), so the two routes give different answers.}
\needspace{14\baselineskip}
\item Survival is 0.8, 0.6 and 0.4 and fecundity 2, 3 and 5 for AA, Aa and aa. Hardy-Weinberg, \(q\) = 0.5. What is \(\Delta q\) for a?
\nopagebreak\par\noindent{\footnotesize\itshape Topic: General equation and mean fitness}
\par\smallskip\noindent\begin{tabularx}{\linewidth}{@{}YYYY@{}}
\makebox[1.9em][l]{(A)}+0.115 & \makebox[1.9em][l]{\textbf{(B)}}\textbf{+0.028}\ \(\checkmark\) & \makebox[1.9em][l]{(C)}\(-\)0.083 & \makebox[1.9em][l]{(D)}0 \\
\end{tabularx}
\keyans{Fitness = 1.6, 1.8 and 2.0. \(\overline{W}\) = 0.4 + 0.9 + 0.5 = 1.8. \(q'\) = (0.45 + 0.50)/1.8 = 0.528, \(\Delta q\) = +0.028. Survival alone gives \(-\)0.083 and fecundity alone +0.115. Survival favors A but fecundity favors a more.}
\needspace{14\baselineskip}
\item Problem Set 12 form. AA = 1.1, Aa = 1, aa = 1.5. \(q\) = 0.25 for a. What happens?
\nopagebreak\par\noindent{\footnotesize\itshape Topic: Heterozygote disadvantage and stability}
\par\smallskip\noindent\begin{tabularx}{\linewidth}{@{}YY@{}}
\makebox[1.9em][l]{(A)}a settles at 0.25 & \makebox[1.9em][l]{(B)}a falls to 0 \\[3pt]
\makebox[1.9em][l]{(C)}a settles at 0.167 & \makebox[1.9em][l]{\textbf{(D)}}\textbf{a rises to fixation}\ \(\checkmark\) \\
\end{tabularx}
\keyans{\(s_1\) = 0.1 and \(s_2\) = 0.5, so the unstable point is 0.1/0.6 = 0.167. At 0.25, a is above it. An a copy experiences 0.25(1.5) + 0.75(1) = 1.125, and an A copy experiences 0.75(1.1) + 0.25(1) = 1.075, so a gains.}
\needspace{14\baselineskip}
\item Absolute fitnesses are 6, 8 and 2 for AA, Aa and aa. What is the equilibrium frequency of A?
\nopagebreak\par\noindent{\footnotesize\itshape Topic: Heterozygote advantage and \(\hat{q}\)}
\par\smallskip\noindent\begin{tabularx}{\linewidth}{@{}YYYY@{}}
\makebox[1.9em][l]{(A)}0.50 & \makebox[1.9em][l]{(B)}0.33 & \makebox[1.9em][l]{(C)}0.25 & \makebox[1.9em][l]{\textbf{(D)}}\textbf{0.75}\ \(\checkmark\) \\
\end{tabularx}
\keyans{Divide by 8. \(W\) = 0.75, 1, 0.25, so \(s_1\) = 0.25 and \(s_2\) = 0.75. \(\hat{q}\) for a = 0.25/1.0 = 0.25, so \(\hat{p}\) for A = 0.75. Logic check, aa is hurt much more, so a is the rare allele.}
\needspace{14\baselineskip}
\item A population is exactly at its heterozygote-advantage equilibrium. What are the average excesses of the two alleles?
\nopagebreak\par\noindent{\footnotesize\itshape Topic: Average excess}
\par\smallskip\noindent\begin{tabularx}{\linewidth}{@{}YY@{}}
\makebox[1.9em][l]{\textbf{(A)}}\textbf{Both zero}\ \(\checkmark\) & \makebox[1.9em][l]{(B)}Equal and opposite but not zero \\[3pt]
\makebox[1.9em][l]{(C)}Each equals its \(s\) & \makebox[1.9em][l]{(D)}Both positive \\
\end{tabularx}
\keyans{At equilibrium neither allele is changing, and \(\Delta q\) = (\(q\)/\(\overline{W}\))\(\cdot \alpha_{a}\), so \(\alpha_{a}\) = 0. Since \(p \cdot \alpha_{A}\) + \(q \cdot \alpha_{a}\) = 0 always, \(\alpha_{A}\) is zero too.}
\needspace{14\baselineskip}
\item In Gabon \(s_2\) = 1 for HbS/HbS and the observed HbS frequency is 0.105. What \(s_1\) would make 0.105 the exact equilibrium?
\nopagebreak\par\noindent{\footnotesize\itshape Topic: Heterozygote advantage and \(\hat{q}\)}
\par\smallskip\noindent\begin{tabularx}{\linewidth}{@{}YYYY@{}}
\makebox[1.9em][l]{(A)}0.121 & \makebox[1.9em][l]{(B)}0.105 & \makebox[1.9em][l]{(C)}0.895 & \makebox[1.9em][l]{\textbf{(D)}}\textbf{0.117}\ \(\checkmark\) \\
\end{tabularx}
\keyans{\(\hat{q}\) = \(s_1\)/(\(s_1\) + 1) = 0.105, so \(s_1\) = 0.105/0.895 = 0.117. The measured 0.121 is close, which is why selection is the main force at this locus.}
\needspace{14\baselineskip}
\item CCR5. \(q\) = 0.90 for WT, \(s\) = 0.16, and \(\Delta q\) = \(-\)0.00171. What is the frequency of \(\Delta\)32 next generation?
\nopagebreak\par\noindent{\footnotesize\itshape Topic: Human cases and Darwin}
\par\smallskip\noindent\begin{tabularx}{\linewidth}{@{}YYYY@{}}
\makebox[1.9em][l]{\textbf{(A)}}\textbf{0.1017}\ \(\checkmark\) & \makebox[1.9em][l]{(B)}0.1171 & \makebox[1.9em][l]{(C)}0.8983 & \makebox[1.9em][l]{(D)}0.0983 \\
\end{tabularx}
\keyans{\(q'\) = 0.90 \(-\) 0.00171 = 0.89829, so \(p'\) = 0.10171. \(\Delta\)32 rises, slowly, because the favored \(\Delta\)32/\(\Delta\)32 homozygote is only \(p^2\) = 1 percent of people.}
\needspace{14\baselineskip}
\item Brown beetles are favored and aa is brown. Green AA and Aa have \(W\) = 0.78. Which statement is right?
\nopagebreak\par\noindent{\footnotesize\itshape Topic: Special models, \(p\) and \(q\)}
\par\smallskip
\begin{enumerate}[label={}, leftmargin=2.4em, labelwidth=0pt, itemsep=2pt, topsep=0pt]
\item \hspace*{-2.4em}\makebox[2.4em][l]{(A)}\(q\) = freq(a), because a is favored
\item \hspace*{-2.4em}\makebox[2.4em][l]{\textbf{(B)}}\textbf{\(q\) = freq(A), and every special \(\Delta q\) for A is negative}\ \(\checkmark\)
\item \hspace*{-2.4em}\makebox[2.4em][l]{(C)}This is heterozygote advantage
\item \hspace*{-2.4em}\makebox[2.4em][l]{(D)}\(q\) = freq(a), because a is recessive
\end{enumerate}
\keyans{A is the allele selected against, and it is dominant. The course rule makes it \(q\). Lecture 24 slide 8 does exactly this.}
\needspace{14\baselineskip}
\item In a hybrid zone, hybrids between two forms survive worse than either pure form, and a single gene controls the difference. Which description fits best?
\nopagebreak\par\noindent{\footnotesize\itshape Topic: Types of selection}
\par\smallskip
\begin{enumerate}[label={}, leftmargin=2.4em, labelwidth=0pt, itemsep=2pt, topsep=0pt]
\item \hspace*{-2.4em}\makebox[2.4em][l]{(A)}Stabilizing selection
\item \hspace*{-2.4em}\makebox[2.4em][l]{(B)}Heterozygote advantage
\item \hspace*{-2.4em}\makebox[2.4em][l]{\textbf{(C)}}\textbf{Heterozygote disadvantage, which drives one allele to fixation within each area}\ \(\checkmark\)
\item \hspace*{-2.4em}\makebox[2.4em][l]{(D)}Purifying selection on both alleles
\end{enumerate}
\keyans{The heterozygote is the least fit. That is an unstable equilibrium, so each local population tends to fix whichever allele is already common there.}
\needspace{14\baselineskip}
\item An inbred population has \(F\) = 0.5, \(q\) = 0.2, and aa is lethal. What is \(q\) after selection?
\nopagebreak\par\noindent{\footnotesize\itshape Topic: Counts or Hardy-Weinberg}
\par\smallskip\noindent\begin{tabularx}{\linewidth}{@{}YYYY@{}}
\makebox[1.9em][l]{(A)}0.167 & \makebox[1.9em][l]{(B)}0.200 & \makebox[1.9em][l]{\textbf{(C)}}\textbf{0.091}\ \(\checkmark\) & \makebox[1.9em][l]{(D)}0.136 \\
\end{tabularx}
\keyans{With inbreeding, aa = \(q^2\) + Fpq = 0.04 + 0.08 = 0.12, Aa = \(2pq\)(1 \(-\) \(F\)) = 0.16, AA = 0.72. \(\overline{W}\) = 0.88. \(q'\) = \(\tfrac{1}{2}\)(0.16)/0.88 = 0.091. Under Hardy-Weinberg it would be 0.167. Inbreeding exposes more a copies in homozygotes, so selection removes them faster.}
\needspace{14\baselineskip}
\item Two populations just merged, giving 0.45 BB, 0.30 Bb and 0.25 bb, and bb is lethal. What is \(q'\)?
\nopagebreak\par\noindent{\footnotesize\itshape Topic: Counts or Hardy-Weinberg}
\par\smallskip\noindent\begin{tabularx}{\linewidth}{@{}YYYY@{}}
\makebox[1.9em][l]{\textbf{(A)}}\textbf{0.200}\ \(\checkmark\) & \makebox[1.9em][l]{(B)}0.286 & \makebox[1.9em][l]{(C)}0.400 & \makebox[1.9em][l]{(D)}0.150 \\
\end{tabularx}
\keyans{Use the counted frequencies, because a freshly mixed population is not in Hardy-Weinberg. \(q\) = 0.25 + 0.15 = 0.40. \(\overline{W}\) = 0.45 + 0.30 = 0.75. \(q'\) = \(\tfrac{1}{2}\)(0.30)/0.75 = 0.200. The 0.286 is the lethal-recessive shortcut \(q\)/(1 + \(q\)), which assumes Hardy-Weinberg. The 0.150 forgets to divide by \(\overline{W}\).}
\needspace{14\baselineskip}
\item Why does \(\hat{q}\) = \(s_1\)/(\(s_1\) + \(s_2\)) put the coefficient of the other allele's homozygote on top?
\nopagebreak\par\noindent{\footnotesize\itshape Topic: Heterozygote advantage and \(\hat{q}\)}
\par\smallskip
\begin{enumerate}[label={}, leftmargin=2.4em, labelwidth=0pt, itemsep=2pt, topsep=0pt]
\item \hspace*{-2.4em}\makebox[2.4em][l]{\textbf{(A)}}\textbf{At equilibrium \(p \cdot s_1\) = \(q \cdot s_2\), so the allele with the more harmful homozygote must be rarer}\ \(\checkmark\)
\item \hspace*{-2.4em}\makebox[2.4em][l]{(B)}It is a convention with no reason
\item \hspace*{-2.4em}\makebox[2.4em][l]{(C)}Because \(s_1\) is always larger
\item \hspace*{-2.4em}\makebox[2.4em][l]{(D)}Because the heterozygote has \(s\) = 0
\end{enumerate}
\keyans{Setting the two alleles' average fitnesses equal gives \(p \cdot s_1\) = \(q \cdot s_2\). Each allele loses the same amount through its own homozygote, so a harsher homozygote must be visited less often.}
\needspace{14\baselineskip}
\item Which fact makes mean fitness ``after selection'' on Lecture 23 slide 12 legitimate to compute with \(p'^{2}\), \(2p'q'\) and \(q'^{2}\)?
\nopagebreak\par\noindent{\footnotesize\itshape Topic: General equation and mean fitness}
\par\smallskip
\begin{enumerate}[label={}, leftmargin=2.4em, labelwidth=0pt, itemsep=2pt, topsep=0pt]
\item \hspace*{-2.4em}\makebox[2.4em][l]{(A)}The slide used observed frequencies
\item \hspace*{-2.4em}\makebox[2.4em][l]{(B)}\(p'\) and \(q'\) always add to 1
\item \hspace*{-2.4em}\makebox[2.4em][l]{(C)}Selection never changes genotype frequencies
\item \hspace*{-2.4em}\makebox[2.4em][l]{\textbf{(D)}}\textbf{It is the next generation after the survivors mate at random}\ \(\checkmark\)
\end{enumerate}
\keyans{Right after selection the survivors are not in Hardy-Weinberg proportions. One round of random mating restores them, and that next generation's \(\overline{W}\) is what the slide computes.}
\needspace{14\baselineskip}
\item Lecture 24 slide 5 says the average excess model works for all single-locus cases. What does it still require?
\nopagebreak\par\noindent{\footnotesize\itshape Topic: Average excess}
\par\smallskip
\begin{enumerate}[label={}, leftmargin=2.4em, labelwidth=0pt, itemsep=2pt, topsep=0pt]
\item \hspace*{-2.4em}\makebox[2.4em][l]{(A)}Relative fitnesses only
\item \hspace*{-2.4em}\makebox[2.4em][l]{\textbf{(B)}}\textbf{Hardy-Weinberg genotype frequencies}\ \(\checkmark\)
\item \hspace*{-2.4em}\makebox[2.4em][l]{(C)}Heterozygote advantage
\item \hspace*{-2.4em}\makebox[2.4em][l]{(D)}Complete dominance
\end{enumerate}
\keyans{``All single-locus cases'' means any fitness pattern at one gene. The derivation still assumes an allele's partner is drawn at random, with probabilities \(q\) and \(p\).}
\needspace{14\baselineskip}
\item A seed-size trait shows high fitness at the large end, falling steadily toward the small end. A large-seed allele is rising. Name both levels.
\nopagebreak\par\noindent{\footnotesize\itshape Topic: Types of selection}
\par\smallskip
\begin{enumerate}[label={}, leftmargin=2.4em, labelwidth=0pt, itemsep=2pt, topsep=0pt]
\item \hspace*{-2.4em}\makebox[2.4em][l]{\textbf{(A)}}\textbf{Directional on the trait, positive on the allele}\ \(\checkmark\)
\item \hspace*{-2.4em}\makebox[2.4em][l]{(B)}Stabilizing on the trait, balancing on the allele
\item \hspace*{-2.4em}\makebox[2.4em][l]{(C)}Disruptive on the trait, positive on the allele
\item \hspace*{-2.4em}\makebox[2.4em][l]{(D)}Directional on the trait, balancing on the allele
\end{enumerate}
\keyans{One extreme favored is directional. A favored allele rising is positive. The two names describe the same event at two levels.}
\needspace{14\baselineskip}
\item A student says selection has stopped in humans because of medicine. What is the best correction?
\nopagebreak\par\noindent{\footnotesize\itshape Topic: Human cases and Darwin}
\par\smallskip
\begin{enumerate}[label={}, leftmargin=2.4em, labelwidth=0pt, itemsep=2pt, topsep=0pt]
\item \hspace*{-2.4em}\makebox[2.4em][l]{(A)}Selection only acts in Africa
\item \hspace*{-2.4em}\makebox[2.4em][l]{(B)}Only mutation acts on humans now
\item \hspace*{-2.4em}\makebox[2.4em][l]{\textbf{(C)}}\textbf{Medicine changes which traits affect fitness, but any heritable trait that still affects survival or reproduction is under selection}\ \(\checkmark\)
\item \hspace*{-2.4em}\makebox[2.4em][l]{(D)}The student is right
\end{enumerate}
\keyans{Elguero and colleagues measured malaria selecting for sickle cell trait in living people in 2015. The three requirements are still met for many traits.}
\end{enumerate}

\newpage
\begin{center}
  {\Large\bfseries Formula Sheet}\\[3pt]
  {\small\itshape Use this page only if your exam provides a formula sheet.}\\[2pt]
  \rule{\textwidth}{0.4pt}
\end{center}

\subsection*{Fitness}
\begin{center}\renewcommand{\arraystretch}{2.1}
\begin{tabular}{r@{\;}l@{\qquad}l}
  \(w\) & \(=\ \text{viability} \times \text{fecundity}\) & absolute fitness, offspring left per zygote\\
  \(W\) & \(=\ w / w_{\max}\) & relative fitness\\
  \(s\) & \(=\ 1 - W\) & selection coefficient
\end{tabular}\end{center}

\subsection*{Mean fitness and the general equation}
\begin{align*}
  \overline{W} &= f(AA)\,W_{AA} + f(Aa)\,W_{Aa} + f(aa)\,W_{aa}\\[4pt]
  q' &= \frac{\tfrac{1}{2}\,f(Aa)\,W_{Aa} + f(aa)\,W_{aa}}{\overline{W}}, \qquad \Delta q = q' - q
\end{align*}
Under Hardy--Weinberg, \(f(AA)=p^2\), \(f(Aa)=2pq\), \(f(aa)=q^2\). With counted individuals that selection acts on directly, use the counted frequencies.

\subsection*{Average excess}
\begin{align*}
  \alpha_a &= q\,W_{aa} + p\,W_{Aa} - \overline{W}, & \Delta q &= \frac{q}{\overline{W}}\,\alpha_a\\[4pt]
  \alpha_A &= p\,W_{AA} + q\,W_{Aa} - \overline{W}, & \Delta p &= \frac{p}{\overline{W}}\,\alpha_A
\end{align*}
Requires Hardy--Weinberg genotype frequencies.

\subsection*{Special models (\(q\) is always the allele selected against)}
\begin{align*}
  \text{Against the recessive:}\quad \Delta q &= \frac{-s\,p\,q^2}{1 - s\,q^2}\\[4pt]
  \text{Against the dominant:}\quad \Delta q &= \frac{-s\,p^2 q}{1 - 2sq + sq^2}
\end{align*}

\subsection*{Heterozygote advantage or disadvantage}
\[
  \hat{q} = \frac{s_1}{s_1 + s_2}
\]
with \(s_1\) for the \(p\) homozygote and \(s_2\) for the \(q\) homozygote. The equilibrium is stable under heterozygote advantage and unstable under heterozygote disadvantage.

\label{LastPage}
\end{document}